% =========================================================
% MICRO-experience — Results
% =========================================================
%
% TODO – Optical microscope images
%   - Include the microscope images provided by Erik (ask him if not yet received)
%   - Caption each image: magnification, scale bar, what structure is shown
%   - Structures visible: parallel trenches at different pitches, circular/radial patterns,
%     alignment marks, contact pads/wiring patterns
%
% TODO – Nominal vs. measured trench width
%   - Table: list nominal trench sizes from the mask (e.g., 2 µm, 3 µm, 5 µm — check
%     the .dat files in piquet/ for measured values)
%   - Measured values from optical microscope or profilometer
%   - Compute the CD (critical dimension) bias = measured − nominal
%   - Discuss origin: optical proximity effects, resist contrast, development time
%
% TODO – Profilometer results
%   - Step height measurement of the photoresist pattern (~1.3 µm expected)
%   - Include profile scan plot (height vs. lateral position)
%   - Correlate with optical image: higher areas = resist remaining on wafer,
%     lower areas = developed/exposed regions (trenches)
%   - Comment on edge roughness / sidewall slope if visible
%
% TODO – Discussion points required by the slides
%   1. Detailed description of the entire photolithography process (refer to M&M)
%   2. Differences between positive-tone and image-reversal process:
%      - undercut profile in image-reversal → useful for lift-off
%      - tone inversion: what was opaque on mask becomes resist, and vice versa
%   3. Nominal vs. final trench size discrepancy — causes and magnitude
%   4. Correlation between optical microscope and profilometer acquisitions:
%      bright (yellow-green) areas in microscope = resist (elevated), dark = substrate (trench)
%
% NOTE: The sample was stored after this experience for use in the DRIE (3rd experience).
